% ============================================================
%  GUÍA 7: TALLER INTEGRADOR Y EVALUACIÓN
%  Curso Introductorio — 3er Año
%  Autor: Ing. Gabriel Astudillo
%  Sesiones cubiertas: 23, 24
% ============================================================

\documentclass[12pt, letterpaper]{article}

% ── Paquetes ─────────────────────────────────────────────────

\usepackage{mathwizards-guia}
\mwguia{Guía 7 — Taller Integrador y Evaluación}{Ing. Gabriel Astudillo $\cdot$ Curso 3er Año}

\begin{document}
% ============================================================

% ── Portada / Título ─────────────────────────────────────────
\mwportada{Guía de Estudio 7}{Taller Integrador y Evaluación Final}{Problemas Complejos que Entrelazan Todas las Disciplinas}

\vspace{4pt}
\begin{cuadroinfo}
\small
\textbf{Presentación:} Esta es la guía final del curso. Aquí no hay temas nuevos: lo que hay son \textbf{problemas que combinan todo lo aprendido}. Densidad con conversión de unidades, cinemática con despeje de fórmulas, lógica con análisis de situaciones científicas. El objetivo es que demuestres que puedes integrar el pensamiento lógico, la manipulación algebraica y el conocimiento científico para resolver problemas complejos usando el Método de Pólya.
\\[4pt]
\textbf{Instrucción:} Para \textbf{cada problema} de esta guía, aplica explícitamente las cuatro fases del Método de Pólya: \textbf{Comprender}, \textbf{Planificar}, \textbf{Ejecutar} y \textbf{Verificar}.
\end{cuadroinfo}

\vspace{8pt}

% ============================================================
\section{Taller Avanzado de Resolución de Problemas}
% ============================================================

\begin{cuadroadvertencia}
\textbf{Recordatorio — El Método de Pólya:}
\begin{enumerate}[label=\textbf{Fase \arabic*:}, leftmargin=*]
  \item \textbf{Comprender:} \textquestiondown Qué datos tengo? \textquestiondown Qué me piden? \textquestiondown Qué unidades necesito?
  \item \textbf{Planificar:} \textquestiondown Qué fórmulas aplico? \textquestiondown Necesito convertir unidades? \textquestiondown Necesito despejar?
  \item \textbf{Ejecutar:} Desarrollo paso a paso, con unidades en cada línea.
  \item \textbf{Verificar:} \textquestiondown El resultado tiene sentido? \textquestiondown Las unidades son coherentes? \textquestiondown Puedo comprobarlo de otra forma?
\end{enumerate}
\end{cuadroadvertencia}

\subsection{Bloque I: Densidad + Conversiones + Geometría}

\begin{enumerate}[leftmargin=*]
  \item \medio{} Un cilindro de aluminio ($d = 2{,}7\;\text{g/cm}^3$) tiene un radio de $5\;\text{cm}$ y una altura de $20\;\text{cm}$.
  \begin{enumerate}[label=\alph*)]
    \item Calcula su volumen en cm$^3$ y en litros.
    \item Calcula su masa en gramos y en kilogramos.
    \item Si lo sumerges completamente en un recipiente con agua, \textquestiondown cuántos mililitros de agua desplaza?
  \end{enumerate}

  \item \medio{} Una esfera metálica tiene una masa de $2{,}5\;\text{kg}$ y una densidad de $7{,}87\;\text{g/cm}^3$.
  \begin{enumerate}[label=\alph*)]
    \item Convierte la masa a gramos.
    \item Calcula el volumen en cm$^3$.
    \item Calcula el radio de la esfera.
    \item \textquestiondown De qué metal podría tratarse?
  \end{enumerate}

  \item \dificil{} En un laboratorio se tiene un bloque de un material desconocido cuyas dimensiones son $4\;\text{in} \times 3\;\text{in} \times 2\;\text{in}$ (pulgadas) y una masa de $1{,}2\;\text{lb}$ (libras). Sabiendo que $1\;\text{in} = 2{,}54\;\text{cm}$ y $1\;\text{lb} = 453{,}6\;\text{g}$:
  \begin{enumerate}[label=\alph*)]
    \item Convierte las dimensiones a centímetros.
    \item Calcula el volumen en cm$^3$.
    \item Convierte la masa a gramos.
    \item Calcula la densidad en g/cm$^3$ e identifica el posible material.
  \end{enumerate}

  \item \dificil{} Se quiere fabricar una esfera de oro ($d = 19{,}3\;\text{g/cm}^3$) que pese exactamente $1\;\text{kg}$.
  \begin{enumerate}[label=\alph*)]
    \item \textquestiondown Qué volumen debe tener la esfera?
    \item \textquestiondown Cuál debe ser su radio?
    \item Si el oro cuesta $\$\,60\;\text{USD/g}$, \textquestiondown cuánto costaría esta esfera?
  \end{enumerate}
\end{enumerate}

\subsection{Bloque II: Cinemática + Conversiones + Despeje}

\begin{enumerate}[leftmargin=*, resume]
  \item \medio{} Un auto viaja a $120\;\text{km/h}$.
  \begin{enumerate}[label=\alph*)]
    \item Convierte su velocidad a m/s.
    \item Si parte de la posición $x_0 = 500\;\text{m}$, \textquestiondown cuál será su posición después de $2\;\text{min}$?
    \item \textquestiondown Cuánto tiempo (en segundos y en minutos) tardará en llegar a la posición $x = 10\;\text{km}$?
  \end{enumerate}

  \item \medio{} La velocidad del sonido en el agua es $1\,480\;\text{m/s}$ y en el aire es $340\;\text{m/s}$.
  \begin{enumerate}[label=\alph*)]
    \item Expresa ambas velocidades en km/h.
    \item \textquestiondown Cuántas veces más rápido viaja el sonido en el agua que en el aire?
    \item Si un buzo bajo el agua escucha una explosión que ocurrió a $3\;\text{km}$ de distancia, \textquestiondown cuántos segundos tardó el sonido en llegarle?
  \end{enumerate}

  \item \dificil{} Un tren de carga viaja a $45\;\text{mph}$ (millas por hora). Sabiendo que $1\;\text{mi} = 1{,}609\;\text{km}$:
  \begin{enumerate}[label=\alph*)]
    \item Convierte su velocidad a km/h y luego a m/s.
    \item Si el tren debe recorrer $200\;\text{km}$, \textquestiondown cuántas horas y minutos tardará?
    \item Si el tren parte de la posición $x_0 = 0$ a las 6:00\,a.m., \textquestiondown a qué hora pasará por el kilómetro 100?
  \end{enumerate}

  \item \dificil{} Un corredor olímpico recorre $100\;\text{m}$ en $9{,}58\;\text{s}$ (récord mundial de Usain Bolt). Un guepardo puede correr a $120\;\text{km/h}$.
  \begin{enumerate}[label=\alph*)]
    \item Calcula la velocidad media de Bolt en m/s y en km/h.
    \item Convierte la velocidad del guepardo a m/s.
    \item Si Bolt y el guepardo partieran al mismo tiempo, \textquestiondown cuántos metros de ventaja le sacaría el guepardo a Bolt al finalizar los $100\;\text{m}$?
  \end{enumerate}
\end{enumerate}

\subsection{Bloque III: Problemas Interdisciplinarios}

\begin{enumerate}[leftmargin=*, resume]
  \item \medio{} Un químico tiene $500\;\text{mL}$ de una sustancia líquida de densidad desconocida. La coloca en una balanza y obtiene una masa de $650\;\text{g}$.
  \begin{enumerate}[label=\alph*)]
    \item Calcula la densidad de la sustancia.
    \item \textquestiondown Esta sustancia flotaría en el agua? \textquestiondown Y en el mercurio ($d = 13{,}6\;\text{g/cm}^3$)?
    \item Usando la tabla: $d_{\text{glicerina}} = 1{,}26\;\text{g/cm}^3$, $d_{\text{aceite}} = 0{,}92\;\text{g/cm}^3$, $d_{\text{ácido sulfúrico}} = 1{,}84\;\text{g/cm}^3$. \textquestiondown De qué sustancia podría tratarse?
  \end{enumerate}

  \item \dificil{} Un estudiante diseña un experimento para medir la velocidad del sonido. Desde un punto A, golpea una campana. Un compañero en el punto B, a $680\;\text{m}$ de distancia, cronometra $2\;\text{s}$ desde que ve el golpe hasta que escucha el sonido.
  \begin{enumerate}[label=\alph*)]
    \item Calcula la velocidad del sonido medida en este experimento.
    \item Convierte el resultado a km/h.
    \item \textquestiondown Por qué se puede ignorar el tiempo que tarda la luz en viajar de A a B? (Pista: calcula ese tiempo sabiendo que la luz viaja a $3 \times 10^8\;\text{m/s}$.)
    \item El valor aceptado es $340\;\text{m/s}$. Calcula el porcentaje de error del experimento.
  \end{enumerate}

  \item \dificil{} Un tanque esférico de radio $0{,}5\;\text{m}$ se llena completamente con un líquido cuya densidad es $0{,}85\;\text{kg/L}$.
  \begin{enumerate}[label=\alph*)]
    \item Calcula el volumen del tanque en m$^3$ y luego conviértelo a litros.
    \item Calcula la masa total del líquido en kilogramos.
    \item Si un camión puede transportar como máximo $3\,000\;\text{kg}$, \textquestiondown cuántos tanques llenos puede llevar?
    \item Expresa la masa de un tanque lleno en libras ($1\;\text{kg} = 2{,}205\;\text{lb}$).
  \end{enumerate}

  \item \dificil{} Un laboratorio de química necesita preparar una experiencia. Tienen una muestra metálica que:
  \begin{itemize}
    \item Pesa $392\;\text{N}$ en la Tierra ($g = 9{,}8\;\text{m/s}^2$).
    \item Al sumergirla en una probeta con $2\;\text{L}$ de agua, el nivel sube a $7\;\text{L}$.
  \end{itemize}
  \begin{enumerate}[label=\alph*)]
    \item Calcula la masa de la muestra.
    \item Calcula su volumen (en litros y en cm$^3$).
    \item Calcula su densidad y determina si podría ser hierro, cobre u oro.
    \item Si la muestra se cortara exactamente por la mitad, \textquestiondown cambiarían su masa, volumen y densidad? Identifica cuáles son extensivas y cuál intensiva.
  \end{enumerate}
\end{enumerate}

\newpage

% ============================================================
\section{Modelo de Prueba Sumativa}
% ============================================================

\begin{cuadroextra}
\textbf{Instrucciones:} Esta sección simula una prueba final del curso. Responde cada pregunta de forma clara, mostrando todos los procedimientos. Tiempo sugerido: 90 minutos.
\end{cuadroextra}

\subsection{Parte I: Preguntas Conceptuales (2 puntos c/u)}

\begin{enumerate}[leftmargin=*]
  \item Define \textbf{proposición} en lógica y da un ejemplo de un enunciado que NO sea proposición. Justifica.

  \item Enumera las cuatro fases del Método de Pólya.

  \item \textquestiondown Cuál es la diferencia entre \textbf{distancia recorrida} y \textbf{desplazamiento}?

  \item Clasifica como propiedad intensiva o extensiva: masa, densidad, temperatura de fusión, volumen.

  \item \textquestiondown Cuál es la diferencia entre una mezcla homogénea y una heterogénea? Da un ejemplo de cada una.

  \item Explica la Tercera Ley de Newton con un ejemplo cotidiano.

  \item \textquestiondown Cuál es la diferencia entre calor y temperatura?

  \item \textquestiondown Qué es una transformación química? Menciona dos indicios que la delaten.
\end{enumerate}

\subsection{Parte II: Problemas de Desarrollo (5 puntos c/u)}

\begin{enumerate}[leftmargin=*, resume]
  \item Resuelve la ecuación: $\dfrac{3x - 1}{2} = \dfrac{x + 5}{4}$

  \item Simplifica usando las propiedades de los exponentes y expresa en notación científica:
  $$\frac{(3 \times 10^5)^2 \times (2 \times 10^{-3})}{6 \times 10^4}$$

  \item Despeja $v_0$ de la fórmula $x = v_0 t + \frac{1}{2}at^2$.

  \item Convierte $250\;\text{km/h}$ a m/s. Luego calcula cuánto recorre un auto a esa velocidad en $45\;\text{s}$.

  \item Un cubo metálico tiene $4\;\text{cm}$ de lado y una masa de $172{,}5\;\text{g}$.
  \begin{enumerate}[label=\alph*)]
    \item Calcula su volumen.
    \item Calcula su densidad.
    \item \textquestiondown Flotaría en mercurio ($d = 13{,}6\;\text{g/cm}^3$)?
  \end{enumerate}

  \item Resuelve el sistema de ecuaciones:
  $\begin{cases} 2x + 3y = 19 \\ 4x - y = 3 \end{cases}$

  \item Un tren parte de una estación a las 7:30\,a.m. a una velocidad constante de $80\;\text{km/h}$. La siguiente estación está a $340\;\text{km}$. \textquestiondown A qué hora llega?

  \item Se mezclan $400\;\text{g}$ de agua a $90\;\text{°C}$ con $600\;\text{g}$ de agua a $20\text{°C}$. Calcula la temperatura de equilibrio. ($c_{\text{agua}} = 1\;\text{cal/(g$\cdot$°C)}$)
\end{enumerate}

\newpage

% ============================================================
\ifmwclaves
\section{Clave de Respuestas}
% ============================================================

\subsection*{Taller Integrador (Problemas 1--12)}

\begin{enumerate}[leftmargin=*, label=\textbf{\arabic*.}]
  \item \begin{enumerate}[label=\alph*)]
    \item $V = \pi r^2 h = \pi(5)^2(20) = 500\pi \approx 1\,570{,}8\;\text{cm}^3 = 1{,}571\;\text{L}$.
    \item $m = d \times V = 2{,}7 \times 1\,570{,}8 \approx 4\,241{,}2\;\text{g} = 4{,}241\;\text{kg}$.
    \item Desplaza $1\,570{,}8\;\text{mL}$ de agua (el volumen del cilindro).
  \end{enumerate}

  \item \begin{enumerate}[label=\alph*)]
    \item $m = 2\,500\;\text{g}$.
    \item $V = m/d = 2\,500/7{,}87 \approx 317{,}7\;\text{cm}^3$.
    \item $V = \frac{4}{3}\pi r^3 \Rightarrow r = \sqrt[3]{\frac{3V}{4\pi}} = \sqrt[3]{\frac{3(317{,}7)}{4\pi}} \approx \sqrt[3]{75{,}8} \approx 4{,}23\;\text{cm}$.
    \item Densidad $7{,}87\;\text{g/cm}^3$ corresponde al hierro.
  \end{enumerate}

  \item \begin{enumerate}[label=\alph*)]
    \item $4 \times 2{,}54 = 10{,}16\;\text{cm}$; $3 \times 2{,}54 = 7{,}62\;\text{cm}$; $2 \times 2{,}54 = 5{,}08\;\text{cm}$.
    \item $V = 10{,}16 \times 7{,}62 \times 5{,}08 \approx 393{,}3\;\text{cm}^3$.
    \item $m = 1{,}2 \times 453{,}6 = 544{,}3\;\text{g}$.
    \item $d = 544{,}3 / 393{,}3 \approx 1{,}38\;\text{g/cm}^3$. Podría ser algún tipo de plástico denso o magnesio.
  \end{enumerate}

  \item \begin{enumerate}[label=\alph*)]
    \item $V = m/d = 1\,000/19{,}3 \approx 51{,}8\;\text{cm}^3$.
    \item $r = \sqrt[3]{\frac{3(51{,}8)}{4\pi}} \approx \sqrt[3]{12{,}37} \approx 2{,}31\;\text{cm}$.
    \item Costo $= 1\,000 \times 60 = \$\,60\,000\;\text{USD}$.
  \end{enumerate}

  \item \begin{enumerate}[label=\alph*)]
    \item $v = 120/3{,}6 = 33{,}\overline{3}\;\text{m/s}$.
    \item $t = 2\;\text{min} = 120\;\text{s}$. $x_f = 500 + 33{,}33 \times 120 = 500 + 4\,000 = 4\,500\;\text{m}$.
    \item $10\,000 = 500 + 33{,}33 \times t \Rightarrow t = 9\,500/33{,}33 = 285\;\text{s} = 4\;\text{min}\;45\;\text{s}$.
  \end{enumerate}

  \item \begin{enumerate}[label=\alph*)]
    \item Agua: $1\,480 \times 3{,}6 = 5\,328\;\text{km/h}$. Aire: $340 \times 3{,}6 = 1\,224\;\text{km/h}$.
    \item $1\,480/340 \approx 4{,}35$ veces.
    \item $t = 3\,000/1\,480 \approx 2{,}03\;\text{s}$.
  \end{enumerate}

  \item \begin{enumerate}[label=\alph*)]
    \item $45 \times 1{,}609 = 72{,}4\;\text{km/h}$. En m/s: $72{,}4/3{,}6 \approx 20{,}1\;\text{m/s}$.
    \item $t = 200/72{,}4 \approx 2{,}76\;\text{h} = 2\;\text{h}\;45\;\text{min}$ (aprox.).
    \item $t_{100} = 100/72{,}4 \approx 1{,}38\;\text{h} \approx 1\;\text{h}\;23\;\text{min}$. Hora: $\approx$ 7:23\,a.m.
  \end{enumerate}

  \item \begin{enumerate}[label=\alph*)]
    \item $v_{\text{Bolt}} = 100/9{,}58 \approx 10{,}44\;\text{m/s} = 37{,}58\;\text{km/h}$.
    \item $v_{\text{guepardo}} = 120/3{,}6 = 33{,}33\;\text{m/s}$.
    \item Bolt tarda $9{,}58\;\text{s}$. En ese tiempo el guepardo recorre: $33{,}33 \times 9{,}58 \approx 319{,}3\;\text{m}$. Ventaja: $319{,}3 - 100 = 219{,}3\;\text{m}$.
  \end{enumerate}

  \item \begin{enumerate}[label=\alph*)]
    \item $d = 650/500 = 1{,}3\;\text{g/cm}^3$.
    \item Se hunde en agua ($d > 1$). Flota en mercurio ($1{,}3 < 13{,}6$).
    \item $d \approx 1{,}3\;\text{g/cm}^3$ se acerca a la glicerina ($1{,}26$).
  \end{enumerate}

  \item \begin{enumerate}[label=\alph*)]
    \item $v = 680/2 = 340\;\text{m/s}$.
    \item $340 \times 3{,}6 = 1\,224\;\text{km/h}$.
    \item $t_{\text{luz}} = 680/(3 \times 10^8) \approx 2{,}27 \times 10^{-6}\;\text{s}$. Es despreciable frente a $2\;\text{s}$.
    \item Error $= \frac{|340 - 340|}{340} \times 100 = 0\%$. El resultado coincide con el valor aceptado.
  \end{enumerate}

  \item \begin{enumerate}[label=\alph*)]
    \item $V = \frac{4}{3}\pi(0{,}5)^3 = \frac{4}{3}\pi(0{,}125) \approx 0{,}5236\;\text{m}^3$. En litros: $0{,}5236 \times 1\,000 = 523{,}6\;\text{L}$.
    \item $m = d \times V = 0{,}85 \times 523{,}6 = 445{,}1\;\text{kg}$.
    \item $3\,000/445{,}1 \approx 6{,}74$, es decir, 6 tanques.
    \item $445{,}1 \times 2{,}205 \approx 981{,}4\;\text{lb}$.
  \end{enumerate}

  \item \begin{enumerate}[label=\alph*)]
    \item $m = W/g = 392/9{,}8 = 40\;\text{kg}$.
    \item $V = 7 - 2 = 5\;\text{L} = 5\,000\;\text{cm}^3$.
    \item $d = 40\,000\;\text{g}/5\,000\;\text{cm}^3 = 8\;\text{g/cm}^3$. Cercano al cobre ($8{,}96$), pero también cercano al hierro ($7{,}87$). Más probablemente hierro o acero.
    \item Masa y volumen se reducen a la mitad (extensivas). La densidad permanece igual (intensiva).
  \end{enumerate}
\end{enumerate}

\subsection*{Modelo de Prueba Sumativa}

\subsubsection*{Parte I: Conceptuales}

\begin{enumerate}[leftmargin=*, label=\textbf{\arabic*.}]
  \item Una proposición es un enunciado declarativo al que se le asigna V o F. «\textquestiondown Qué hora es?» no es proposición porque es una pregunta, no afirma nada.

  \item Comprender, Planificar, Ejecutar, Verificar.

  \item La distancia recorrida es la longitud total del camino (escalar, siempre $\geq 0$). El desplazamiento es el cambio de posición (vectorial, puede ser negativo o cero).

  \item Masa: extensiva. Densidad: intensiva. Temperatura de fusión: intensiva. Volumen: extensiva.

  \item Homogénea: una sola fase, no se distinguen componentes (ej. agua con sal). Heterogénea: dos o más fases visibles (ej. agua con aceite).

  \item Al caminar, tus pies empujan el suelo hacia atrás (acción) y el suelo te empuja hacia adelante (reacción), permitiéndote avanzar.

  \item Temperatura: grado de agitación de las partículas. Calor: transferencia de energía entre cuerpos a diferente temperatura.

  \item Es un cambio que altera la composición de la materia, formando sustancias nuevas. Indicios: cambio de color permanente y producción de gas.
\end{enumerate}

\subsubsection*{Parte II: Problemas de Desarrollo}

\begin{enumerate}[leftmargin=*, label=\textbf{\arabic*.}]
  \setcounter{enumi}{8}
  \item Multiplicamos por 4: $2(3x - 1) = x + 5 \Rightarrow 6x - 2 = x + 5 \Rightarrow 5x = 7 \Rightarrow x = \frac{7}{5} = 1{,}4$.

  \item $\frac{(3)^2 \times (10^5)^2 \times 2 \times 10^{-3}}{6 \times 10^4} = \frac{9 \times 10^{10} \times 2 \times 10^{-3}}{6 \times 10^4} = \frac{18 \times 10^7}{6 \times 10^4} = 3 \times 10^3$.

  \item $x - \frac{1}{2}at^2 = v_0 t \Rightarrow v_0 = \dfrac{x - \frac{1}{2}at^2}{t} = \dfrac{2x - at^2}{2t}$.

  \item $v = 250/3{,}6 \approx 69{,}44\;\text{m/s}$. Distancia $= 69{,}44 \times 45 \approx 3\,125\;\text{m} = 3{,}125\;\text{km}$.

  \item a) $V = 4^3 = 64\;\text{cm}^3$. b) $d = 172{,}5/64 \approx 2{,}7\;\text{g/cm}^3$. c) Sí flotaría ($2{,}7 < 13{,}6$).

  \item De la segunda: $y = 4x - 3$. Sustituimos: $2x + 3(4x - 3) = 19 \Rightarrow 14x - 9 = 19 \Rightarrow 14x = 28 \Rightarrow x = 2$. $y = 4(2) - 3 = 5$.

  \item $t = 340/80 = 4{,}25\;\text{h} = 4\;\text{h}\;15\;\text{min}$. Hora de llegada: 7:30 + 4:15 = 11:45\,a.m.

  \item $Q_{\text{cedido}} = Q_{\text{absorbido}}$: $400(90 - T_f) = 600(T_f - 20)$. $36\,000 - 400T_f = 600T_f - 12\,000$. $48\,000 = 1\,000T_f$. $T_f = 48\;$.
\end{enumerate}

\fi
\end{document}
